%%%%%%%%%%%%%%%%%%%%%%%%%%%%%%%%%
%
%       EXERCÍCIO
%
%%%%%%%%%%%%%%%%%%%%%%%%%%%%%%%%%

\ifdefstring{\atividade}{prova}{%
    \ifdefstring{\modo}{objetivo}{%
        \renewcommand{\valorquestao}{\ValorQObj\ ponto}
    }{%
        \renewcommand{\valorquestao}{\ValorQDisc\ pontos}
    }%
}%

\begin{exercicioBanco}[\valorquestao]
Determine a equação da circunferência de centro \(C = (5,-5)\) e que contém o ponto \(P = (0,-5)\).

% Define as alternativas
\newcommand{\alternativas}{%
    \begin{center}
        \begin{tabularx}{\textwidth}{XXX}
            (a) \(\dfrac{x^{2}}{5} + \dfrac{y^{2}}{5} - 2x + 2y + 5 = 0\). &
            (b) \(x^{2} - 10x + y^{2} + 10y = 25\). &
            (c) \(x^{2} - x + y^{2} - y - 5 = 0\). \\[15pt]
            (d) \(5x^{2} + 5y^{2} - x - y = 25\). &
            (e) \(5x^{2} + 5y^{2} + x + y = 25\).
        \end{tabularx}
    \end{center}
}

% Define a resposta correta
\newcommand{\resposta}{A}

% Lógica condicional para exibição
\ifdefstring{\atividade}{lista}{%
    \alternativas
    \vspace{0.5em}
    
    \noindent\textbf{Resposta:} letra \textbf{\resposta}.
}{%
    \ifdefstring{\modo}{objetivo}{%
        \alternativas
    }{%
        % modo = discursiva → não mostra alternativas
    }
}
\end{exercicioBanco}

